%% 摘要

% 中文摘要
\begin{abstract}
大语言模型通常需要后训练以符合人类偏好和任务需求。现有偏好对齐方法主要包括监督微调、基于人类反馈的强化学习和直接偏好优化等。其中，基于人类反馈的强化学习通常先从人工偏好比较数据中恢复奖励模型，再将其固定用于在线策略优化。这一流程既会继承偏好标注中的主观差异和噪声，也会在策略生成分布逐渐偏离奖励模型训练分布后降低评分可靠性，并进一步引发奖励模型过优化和奖励黑客问题。因此，偏好对齐的难点不仅在于单轮偏好拟合精度，还在于策略持续变化时奖励信号能否保持稳定、可控的更新幅度。

针对上述问题，本文研究大语言模型偏好对齐中的奖励函数学习问题，目标是从专家示范数据中恢复可用于策略优化的奖励信号，并在策略分布持续变化时保持该信号稳定。逆强化学习将奖励函数视为可学习对象，能够在策略更新的同时调整奖励；但若缺少奖励变化约束，相邻轮次奖励函数容易出现尺度膨胀或语义漂移，从而破坏训练稳定性。基于此，本文提出近端逆奖励优化方法PIRO（Proximal Inverse Reward Optimization）。PIRO以专家示范响应作为偏好行为证据，在逆强化学习框架下直接从示范数据中学习奖励函数，减少对额外成对偏好标注的依赖，从而降低偏好比较环节可能引入的主观误差；同时，PIRO限制相邻迭代轮次之间的奖励变化幅度，使奖励模型适应当前策略分布并避免无约束漂移。理论上，本文将对齐过程建模为奖励学习与策略优化耦合的双层问题，推导局部替代目标与原始目标之间的改进下界，说明奖励变化误差项会影响原始对齐目标的稳定改进。算法上，本文用相邻轮次奖励输出的均方根差（RMSD）近似奖励变化幅度，并设计自适应约束系数调节机制，将理论误差控制项转化为训练过程中可观测、可调节的稳定性反馈信号。

本文在TL;DR摘要生成和UltraFeedback通用指令跟随两个任务上进行验证。TL;DR任务中，PIRO在iter5取得奖励评分$1.415$，相较固定奖励RLHF/PPO提升约$14.1\%$，相较SFT提升约$30.0\%$；去除奖励变化约束后，模型早期短暂上升，随后持续退化，说明该约束对多轮IRL训练稳定性具有必要作用。UltraFeedback任务中，PIRO最佳检查点的MT-Bench Avg(8)为$7.53$，超过SFT、DPO、RLHF(PPO)、RLHF(GRPO)以及无约束IRL基线，并在推理和STEM等维度上取得较明显提升。实验结果表明，PIRO能够在动态奖励学习中抑制奖励漂移，提升多轮逆强化学习偏好对齐训练的稳定性。

\keywordslist{大语言模型对齐, 逆强化学习, 奖励模型, RMSD约束, 策略优化}
\end{abstract}

% 英文摘要
\begin{engabstract}
Large language models require post-training to match human preferences and task requirements. Common alignment methods include supervised fine-tuning, reinforcement learning from human feedback, and direct preference optimization. In reinforcement learning from human feedback, a reward model is learned from human preference comparisons and then fixed during policy optimization. The recovered reward inherits annotation subjectivity and noise; as the policy leaves the reward-model training distribution, evaluation becomes less reliable and may lead to reward overoptimization and reward hacking.

Motivated by this problem, this thesis studies reward-function learning for preference alignment of large language models: recovering a reward signal from expert demonstrations and keeping it stable as the policy distribution changes. Inverse reinforcement learning (IRL) can update reward and policy jointly, but unconstrained reward updates may inflate reward scales or shift reward semantics. This thesis therefore proposes Proximal Inverse Reward Optimization (PIRO). PIRO learns the reward function directly from expert demonstration responses, reducing reliance on additional pairwise preference labels and their annotation error. It also bounds reward changes between adjacent iterations, allowing adaptation to the current policy while avoiding unconstrained drift. The analysis formulates alignment as a bilevel problem and derives a local lower bound linking reward-change error to improvement of the original objective. The algorithm uses the root mean square difference (RMSD) between reward outputs and an adaptive coefficient as an observable stability signal.

The method is evaluated on TL;DR summarization with Pythia-1B and UltraFeedback instruction following with Qwen3.5-9B. On TL;DR, PIRO reaches $1.415 \pm 0.054$ at iter5, $14.1\%$ above RLHF/PPO and $30.0\%$ above SFT. Removing the $C\epsilon$ constraint gives an early gain followed by degradation, indicating the necessity of constraining reward changes. On UltraFeedback, PIRO's best checkpoint scores $7.53$ on MT-Bench Avg(8), outperforming SFT, DPO, RLHF(PPO), RLHF(GRPO), and the unconstrained IRL baseline, with clearer gains on reasoning and STEM. These results indicate that PIRO suppresses reward drift and improves multi-round IRL-based preference alignment.

\engkeywordslist{LLM alignment, IRL, reward model, RMSD constraint}
\end{engabstract}
